\myparagraph{Weighted automata} Weighted automata generalize finite
automata with weights on transitions~\cite{Droste:2009:HWA:1667106}. 
They generalize word languages to partial functions from words to
values of a semiring. First introduced by Sch\"utzenberger and Chomsky in the 60s,
they have been studied for long~\cite{Droste:2009:HWA:1667106}, with
applications in natural language and image processing for
instance. More recently, they have found new applications in
computer-aided verification as a measure of system
quality through quantitative properties~\cite{cdh10}, and in
system synthesis, as objectives for quantitative games~\cite{conf/concur/FiliotGR12}. 
In this paper, we consider weighted automata $\calN$ over the semiring
$(\mathbb{Z}\cup  \{ -\infty \}, \max,+)$, and just
call them weighted automata (WA). The value of a run is the sum of the weights
occurring on its transitions, and the value of a word is the maximal value of all
its accepting runs.  Absent transitions have a weight of $-\infty$ and runs of
value $-\infty$ are considered non-accepting. This defines a partial function
denoted $\inter{\calN} : \Sigma^*\rightarrow \mathbb{Z}$ whose domain is denoted
by $\lang{\calN}$. 

\myparagraph{Determinization problem} Most of the good algorithmic properties of
finite automata do not transfer to WA. Notably, the
(quantitative) language inclusion $\inter{\mathcal{A}} \leq \inter{\mathcal{B}}$ is undecidable for
WA~\cite{Krob:1992:EPR} (see also \cite{dav16} and
\cite{conf/atva/AlmagorBK11} for different proofs based on reductions
from the halting problem for two-counter machines). This has triggered 
research on sub-classes or other formalisms for which this problem becomes
decidable~\cite{DBLP:conf/fsttcs/FiliotGR14,conf/concur/ChatterjeeDEHR10}.
This includes the class of deterministic WA (DWA, also known as sequential
WA in the literature), which are the WA whose underlying (unweighted)
automaton is deterministic. Another scenario where it is desirable to
have a DWA is the quantitative synthesis problem,  
undecidable even for unambiguous WA, yet decidable for
DWA~\cite{conf/concur/FiliotGR12}. However, and in contrast with finite unweighted
automata, WA are not determinizable in general. For instance, the function
which outputs the maximal value between the number of $a$'s and the
number of $b$'s in a word $\alpha \in\{a,b\}^*$ is not realizable with a DWA. 
This motivates the determinization problem: given a WA $\calN$, is it
determinizable? I.e. is there a DWA defining the same
(partial) function as $\calN$? 

The determinization problem for computational models is
fundamental in theoretical computer science. For WA in particular, it
is sometimes more natural (and at least exponentially more succinct)
to specify a (non-deterministic) WA, even if some equivalent DWA
exists. If the function is specified in a weighted logic equivalent to
WA, such as weighted MSO~\cite{journals/tcs/DrosteG07}, the
logic-to-automata transformation may construct a non-deterministic,
but determinizable, WA. However, despite many research efforts, the largest class
for which this problem is known to be decidable is the class of
polynomially ambiguous WA~\cite{kirsten_et_al:LIPIcs:2009:1850}, and
the decidability status for the full class of WA is a long-standing
open problem. Other contributions and approaches to the determinization problem
%In addition to the contribution on polynomially
%ambiguous WA mentioned before, there have been other contributions to
%the determinization problem. These 
include the identification of sufficient conditions for
determinizability~\cite{Mohri_FiniteState_CoLing}, approximate determinizability
(for unambiguous WA) where the DWA is required to produce values %that are at
at most $t$ times the value of the WA, for a given factor
$t$~\cite{journals/tcs/AminofKL13}, and (incomplete) approximation algorithms
when the weights are non-negative~\cite{journals/mst/ColcombetD16}.

\myparagraph{Bounded-delay \& regret determinizers} In this paper, we adopt
another approach that consists in constraining the class of DWA that can be
used for determinization. More precisely, we define a \emph{class} of
DWA $\mathfrak{C}$ as a function from WA to sets of DWA, and
say that a DWA $\calD$ is a $\mathfrak{C}$-determinizer of a WA $\calN$ if
\begin{inparaenum}[$(i)$]
\item $\calD \in \mathfrak{C}(\calN)$ and
\item $\inter{\calN} =
\inter{\calD}$.
\end{inparaenum}
Then, $\calN$ is said to be $\mathfrak{C}$-deteminizable if it admits a
${\mathfrak{C}}$-determinizer. If ${\bf DWA}$ denotes the function mapping any
WA to the whole set of DWA, then obviously the ${\bf DWA}$-determinization
problem is the general (open) determinization problem. In this paper, we consider
two restrictions. 

First, given a bound $k\in\mathbb{N}$, we look for the class of $k$-delay DWA
${\bf Del}_k$, which maintain a strong relation with the sequence of values
along some accepting run of the non-deterministic automaton. More
precisely, a DWA $\calD$ belongs to ${\bf Del}_k(\calN)$ if for all words $\alpha \in
\lang{\calD}$, there is an accepting run $\rho$ of $\calN$ with maximal value
such that the running sum of the prefixes of $\rho$ and the running sum of the
prefixes of the unique run $\rho_{\calD}$ of $\calD$ on $\alpha$ are constantly
close in the following sense: for all lengths $\ell$, the absolute value of the
difference of the value of the prefix of $\rho$ of length $\ell$ and
the value of the prefix of $\rho_\calD$ of length $\ell$ is at most $k$.
Then the ${\bf
Del}_k$-determinization problem amounts to deciding whether there exists
${\calD} \in {\bf Del}_k$ such that $\inter{\calN}=\inter{\calD}$.  And if $k$
is left unspecified, it amounts to decide whether there exists ${\calD} \in
\bigcup_{k \in {\mathbb{N}}} {\bf Del}_k(\calN)$ such that
$\inter{\calN}=\inter{\calD}$. We note ${\bf Del}$ the function mapping any WA
$\calN$ to $\bigcup_k {\bf Del}_k(\calN)$. We will show that the ${\bf
Del}_k$-determinization problem is  complete for \EXP, and the ${\bf
Del}$-determinization problem is equivalent to the general ({\it unconstrained})
determinization problem.

The notion of delay has been central in many works on automata with
outputs. For instance, it has been a key notion in transducer theory
(automata with word outputs) for the determinization and the
functionality problems \cite{journals/tcs/BealCPS03},  and the decomposition of finite-valued
transducers~\cite{sakarovitch_et_al:LIPIcs:2008:1324}. The notion of delay has been also used in the
theory of WA, for instance to give sufficient conditions for
determinizability~\cite{Mohri_FiniteState_CoLing} or for the decomposition of finite-valued group
automata~\cite{DBLP:conf/fsttcs/FiliotGR14}. 

\begin{example} Let $A = \{a,b\}$ and
    $k\in\mathbb{N}$. The left automaton of Fig.~\ref{fig:example2}
    maps any word in $AaA^*$ to $0$, and any word in $AbA^*$ to
    $1$. It is ${\bf Del}_k$-determinizable by the right automaton of
    Fig.~\ref{fig:example2}. After one step, the delay of the DWA is
    $k$ with both transitions of the left WA. After two or more steps,
    the delay is always $0$. It is not ${\bf Del}_j$-determinizable for
    any $j<k$.
\begin{figure}[h]
\begin{center}
\begin{tikzpicture}[node distance=0.8cm]
	\node[state,initial left] (si) {};
	\node[state,above right= of si] (s0) {};
	\node[state,below right= of si] (s1) {};
	\node[state,right= of s0,accepting] (s2) {};
	\node[state,right= of s1] (s3) {};

	\path
	(si) edge node[el] {$A,-k$} (s0)
	(si) edge node[el,swap] {$A,k$} (s1)
	(s0) edge node[el] {$a,k$} (s2)
	(s1) edge node[el,swap,pos=0.8] {$b,1-k$} (s2)
	(s0) edge node[el,pos=0.8] {$b,0$} (s3)
	(s1) edge node[el,swap] {$a,0$} (s3)
	(s2) edge [loopright] node[el] {$A,0$} (s2)
	(s3) edge [loopright] node[el] {$A,0$} (s3)
	;

	\node[state,initial above] at (4,0) (qi) {};
	\node[state,right= of qi]  (q) {};
	\node[state,accepting, right= of q] (p) {};
        
        \path
	(qi) edge node[el] {$A,0$} (q)
	(q) edge node[el] {$a,0$} node[el,below] {$b,1$} (p)
	(p) edge [loopabove] node[el] {$A,0$} (p)
        ;
\end{tikzpicture}
\end{center}
%A %complete and
% \caption{A pair-deterministic automaton with alphabet $A =
% \{a,b\}$. For all $k \in \mathbb{N}$, the automaton is $k$-delay
% determinizable but not $k$-regret determinizable.}
%\vspace{-3mm}
\caption{A WA (left) and one of its $k$-delay determinizers (right)}
\label{fig:example2}
\end{figure}
\end{example}


Second, we consider the class $\mathbf{Hom}$ of so-called homomorphic
DWA. Intuitively, any DWA $\calD\in \mathbf{Hom}(\calN)$
maintains a close relation with the structure of $\calN$: the
existence of an homomorphism from $\calD$ to $\calN$. An alternative definition
is that of a $0$-\emph{regret game} \cite{hpr15-journal} played on $\calN$:
\adam chooses input symbols one by one (forming a word $\alpha \in
\lang{\calN}$), while \eve reacts by choosing transitions of $\calN$, thus
constructing a run $\rho$ of ${\calN}$ on the fly (\ie without knowing the full
word $\alpha$ in advance). \eve wins the game if $\rho$ is accepting and its
value is equal to $\inter{\calN}(\alpha)$, \ie $\rho$ is a maximal accepting run
on $\alpha$. Then, any (finite memory) winning strategy for \eve can be seen as
a $\mathbf{Hom}$-determinizer of $\calN$ and conversely.  This generalizes the
notion of good-for-games automata, which do not need to be determinized prior to
being used as observers in a game, from the Boolean setting~\cite{hp06} to the
quantitative one. In some sense, $\mathbf{Hom}$-determinizable WA are ``good for
quantitative games'': when used as an observer in a quantitative game, Eve's
strategy can be applied on the fly instead of determinizing the WA and
constructing the synchronized product of the resulting DWA with the game arena.
This notion has been first introduced in~\cite{akl10} with motivations coming
from the analysis of online algorithms.  In~\cite{akl10},  it was shown that the
$\mathbf{Hom}$-determinization problem is in \P.

\begin{example}\label{ex:regret}
The following WA maps both $ab$ and $aa$ to $0$. It is not 
$\mathbf{Hom}$-determinizable because Eve has to choose whether to go
left or right on reading $a$. If she goes right, then Adam wins by
choosing letter $b$. If she goes left, Adam wins by picking $a$
again. However, it is \emph{almost} $\mathbf{Hom}$-determinizable by the DWA
obtained by removing the right part, in the sense that the function
realized by this DWA is $1$-close from the original one. This
motivates approximate determinization. 
\begin{figure}[h]
\begin{center}
\begin{tikzpicture}[node distance=0.8cm]
	\node[state,initial above] (si) {};
	\node[state,right= of si] (s0) {};
	\node[state,left= of si] (s1) {};
	\node[state,right= of s0,accepting] (s2) {};
	\node[state,left= of s1,accepting] (s3) {};

	\path
	(si) edge node[el] {$a,0$} (s0)
	(si) edge node[el,swap] {$a,0$} (s1)
	(s0) edge node[el] {$a,1$} (s2)
	(s1) edge node[el,swap,align=left,pos=0.3] {$a,0$\\$b,0$} (s3)
	;
\end{tikzpicture}
\end{center}
% \caption{A pair-deterministic automaton with alphabet $A =
% \{a,b\}$. For all $k \in \mathbb{N}$, it is $1$-regret and $k$-delay
% determinizable, yet for all $0 \le j < k$ it is not $j$-delay determinizable.}
%\vspace{-2mm}
\caption{A WA that is $\mathbf{Del}_0$-determinizable and 
$(1,\mathbf{Hom})$-determinizable, but not $\mathbf{Hom}$-determinizable}
\label{fig:example3}
\end{figure}


\end{example}

\myparagraph{Approximate determinization} Approximate
determinization of a WA $\calN$ relaxes the determinization problem
to determinizers which do not define
exactly the same function as $\calN$ but approximate it. Precisely, for a class
$\mathfrak{C}$ of DWA, $\calD$ a DWA and $r\in\mathbb{N}$,
we say that $\calD$ is an
$(r,\mathfrak{C})$-determinizer of $\calN$ if
\begin{inparaenum}[$(i)$]
\item $\calD\in \mathfrak{C}(\calN)$,
\item $\lang{\calD} = \lang{\calN}$ and
\item for all words $\alpha \in \lang{\calN}$, $|\inter{\calN}(\alpha)-\inter{\calD}(\alpha)|\leq r$.
\end{inparaenum}
Then, $\calN$ is $(r,\mathfrak{C})$-determinizable if
it admits some $(r,\mathfrak{C})$-determinizer, and it is 
\emph{approximately} $\mathfrak{C}$-determinizable if it is 
$(r,\mathfrak{C})$-determinizable for some $r$. 


As Example~\ref{ex:regret} shows, there are WA that are approximately
$\mathbf{Hom}$-determinizable but not $\mathbf{Hom}$-determinizable, making
this notion appealing for the class of homomorphic determinizers. However,
there are classes $\mathfrak{C}$ for which a WA is approximately
$\mathfrak{C}$-determinizable if and only if it is
$\mathfrak{C}$-determinizable, making approximate determinization much less
interesting for such classes. This is the case for classes $\mathfrak{C}$ which
are \emph{complete for determinization} (Theorem~\ref{thm:rclose}), in the sense
that any determinizable WA is also $\mathfrak{C}$-determinizable. Obviously, the
class $\mathbf{DWA}$ is complete for determinization, but we show it is also the
case for the class of bounded-delay determinizers $\mathbf{Del}$
(Theorem~\ref{thm:semi-det}).
%
Therefore, we study approximate determinization for the class of homomorphic
determinizers only. We call such determinizers $r$-regret
determinizers, building on the regret game analogy given above. Indeed, a WA
$\calN$ is $(r,\mathbf{Hom})$-determinizable if and only if \eve wins the regret game
previously defined, with the following modified winning condition: the run that she
constructs on the fly must be such that
$|\inter{\calN}(\alpha)-\inter{\calD}(\alpha)|\leq r$, for all words $\alpha \in
\lang{\calN}$ that \adam can play.
%in the course of the game.
We say that a WA $\calN$ is approximately $\mathbf{Hom}$-determinizable if
there exists a $r \in \mathbb{N}$ such that $\calN$ is
$(r,\mathbf{Hom})$-determinizable.


\myparagraph{Contributions} We show that the ${\bf Del}_k$-determinization problem
problem is \EXP-complete, even when $k$ is fixed
(Theorems~\ref{thm:kdelay-in-exp}~and~\ref{thm:kdelay-exp-hard}).
We also show that the class $\mathbf{Del}$ is
complete for determinization, \ie any determinizable WA $\calN$
is $k$-delay determinizable for some $k$ (Theorem~\ref{thm:semi-det}). This shows that
solving the $\mathbf{Del}$-determinization problem would solve the
(open) general determinization problem. This also gives a new
(complete) semi-algorithm for determinization, which consists in testing for the
existence of $k$-delay determinizers for increasing values $k$.  %Furthermore,
We exhibit a family of bounded-delay determinizable WA, for delays
which depend exponentially on the WA. Despite our efforts, exponential delays
are the highest lower bound we have found. Interestingly, finding higher lower
bounds would lead to a better understanding of the determinization
problem, and proving that one of these lower bound is also an upper bound would
immediately give decidability. 
% Central results for the ${\bf Del}_k$-determinization problem are obtained with
% the help of a pumping lemma related to the notion of delay that is interesting
% on its own (Lemma~\ref{}).\todo{Which lemma is that?}
To decide ${\bf Del}_k$-determinization, we provide
a reduction to $\mathbf{Hom}$-determinization (\ie $0$-regret determinization),
which is known to be decidable in polynomial time~\cite{akl10}. 

We show that the approximate $\mathbf{Hom}$-determinization problem is
decidable in exponential time (Theorem~\ref{thm:existential-regret-prob}), a
problem which was left open in \cite{akl10}. This result is based on a
non-trivial extension to the quantitative setting of a game tool proposed by
Kuperberg and Skrzypczak in~\cite{ks15} for Boolean automata. In particular, our
quantitative extension is based on energy games~\cite{bflms08} while parity
games are sufficient for the Boolean case.  If $r$ is given (in
binary) the $(r,\mathbf{Hom})$-determinization problem is shown to be
\EXP-complete
(Theorems~\ref{thm:regret-upper}~and~\ref{thm:regret-lower-bound}). The hardness
holds even if $r$ is given in unary.  In the course of establishing our results,
we also show that every WA $A$ that is approximately
$\mathbf{Hom}$-determinizable is also {\em exactly} determinizable but there
may not be a homomorphism from a deterministic version of the automaton to the
original one (Lemma~\ref{lem:joker-nec-regret} and
Theorem~\ref{thm:win-joker-det}). Hence, the decision procedure for approximate
$\mathbf{Hom}$-determinizability can also be used as an algorithmically
verifiable sufficient condition for determinizability. 

\myparagraph{Other related works}
In transducer theory, a notion close to the notion of $k$-delay
determinizer has been introduced in \cite{fjlw16}, that of $k$-delay
uniformizers of a transducer. A uniformizer of a transducer $T$ is an
(input)-deterministic transducer such that the word-to-word function it defines
(seen as a binary relation) is included in the relation defined by
$T$, and any of its accepting runs should be $k$-delay close from some
accepting run of $T$. While the notion of $k$-delay uniformizer in
transducer theory is close to the notion of $k$-delay determinizer for
WA, the presence of a max operation in WA makes the $k$-delay
determinization problem conceptually harder. 
